\section{Límites de la extrapolación: otros mecanismos, LSTM y Scaling Laws}

La Induction-head story resulta muy atractiva porque conecta la conducta, la estructura matemática y la intervention; cuanto más pulcra es una historia, más hace falta buscar activamente lo que no logra explicar. Al final, el ponente admite de manera explícita otros mechanisms y deja como preguntas abiertas buena parte del cómputo de los modelos grandes.

\subsection{In-context nearest neighbor no abarca todo el meta-learning}

Olah resume la induction como in-context nearest neighbor y, al mismo tiempo, señala que las soft heads pueden aprender functions. Esta descripción sirve para explicar cómo se empareja y se continúa un pattern a partir de examples, pero no cubre necesariamente la gradient-descent-like inner optimization, el Bayesian updating, el multi-step reasoning ni el uso de herramientas.

\lecturefigure{slide-61-in-context-nearest-neighbor.jpg}{La induction como in-context nearest neighbor; las soft heads pueden aprender relaciones funcionales.}{Video oficial de Stanford Online, 00:55:30--00:55:33.}

\lecturefigure{slide-62-alternative-contributors.jpg}{Otros mechanisms pueden contribuir al mismo tiempo al large-model ICL.}{Video oficial de Stanford Online, 00:55:34--00:55:59.}

\begin{warningbox}{Una mechanistic story elegante puede seguir siendo incompleta}
Las induction heads pueden explicar cierto tipo de lookup/copy behavior y, en los small models, explicar el chosen ICL score; los modelos grandes pueden usar a la vez MLP memory, distributed feature binding, other heads, implicit optimization o paths más largos. La conclusión más adecuada es «un mecanismo importante», no «el único mecanismo».
\end{warningbox}

\teachervoice{La present theory que se presenta al cierre de la clase abarca, en esencia, dos tipos de algoritmos: uno es el copying amplio del tipo «si ya apareció cierta palabra, puede volver a aparecer más adelante en una posición adecuada»; el otro es el induction-head nearest neighbor. El ponente deja espacio para «possibly other things».}

\subsection{Por qué el Transformer es más adecuado que la bounded-state LSTM para el lookup en contextos largos}

La attention puede acceder directamente a cualquier position pasada, y el costo de la recuperación depende sobre todo de la context length y de la attention implementation; la LSTM tradicional debe comprimir toda la historia en un recurrent state de dimensión fija, por lo que el lookup exacto a larga distancia tiende a degradarse. Esta diferencia hace que la induction-style retrieval sea más natural en el Transformer, pero no significa que la LSTM sea matemáticamente incapaz de expresar algún tipo de meta-learning, ni que la attention aprenda automáticamente el lookup correcto.

\lecturefigure{slide-63-transformer-vs-lstm.jpg}{ICL en contextos largos: la ruta de recuperación directa de la attention frente al cuello de botella del estado fijo de la LSTM.}{Video oficial de Stanford Online, 00:56:00--00:56:21.}

\begin{knowledgebox}{No hay que presentar la comparación de arquitecturas como una imposibilidad absoluta}
Con precisión finita, hidden state finito y las condiciones reales de entrenamiento, la retrieval fiel a larga distancia es difícil para la LSTM; en teoría, un estado más grande, una memory externa o un entrenamiento especial podrían cambiar la conclusión. El «impossible» de la clase debe entenderse en un contexto de ingeniería: ese tipo de mechanism no surge de forma natural y es difícil de escalar de manera estable.
\end{knowledgebox}

\subsection{¿Explican las induction heads las meta-learning scaling laws?}

Si el ICL se aproxima a un context nearest neighbor, la capability podría variar en función conjunta de la pattern diversity de los training data, la representation quality y la context retrieval, y no solo del parameter count. El ponente plantea que esto podría ayudar a entender las curvas de las scaling laws, pero no ofrece una teoría completa ni una validación.

\lecturefigure{slide-64-meta-learning-scaling-laws.jpg}{Pregunta abierta: ¿puede el induction-head mechanism explicar las meta-learning scaling laws?}{Video oficial de Stanford Online, 00:56:22--00:56:47.}

\teachervoice{La explicación de Olah sobre las scaling laws es claramente más tentative: se trata solo de una línea posible. Las notas no deben convertir una diapositiva de preguntas en una law establecida, y mucho menos sustituir con la induction story el análisis empírico de la data distribution, la optimization y el architecture scaling.}

\subsection{Lo que todavía no se entiende}

En la sesión de preguntas y respuestas, Olah enumera las MLP layers, la arithmetic, la multi-speaker/dialogue representation y una gran cantidad de mecanismos internos de los modelos grandes. La near-complete account de los attention-only toy models es una demostración del método, no el punto final de la investigación; la nonlinear feature superposition, la normalization y los distributed circuits de los modelos grandes siguen siendo difíciles.

\teachervoice{«Hay tanto que todavía no entendemos que es difícil decir cuál es el siguiente paso» es el cierre más honesto de esta clase. La madurez de la interpretabilidad mecanicista debe medirse por circuits reproducibles, resultados de intervención y casos de falla, no por la cantidad de términos.}

\subsection{Resumen de la sección}

Las induction heads son un ICL mechanism importante, con estructura y respaldo de intervention, pero los modelos grandes pueden usar otros algorithms. El Transformer ofrece una interfaz natural para la retrieval en contextos largos, la conexión con las scaling laws sigue siendo una hipótesis abierta, y las MLP y los distributed circuits son la dificultad central de la siguiente etapa.

